% ============================================================
%  Rapport de Stage de Fin d'Études
%  Rayen Thabet -- Tunisie Telecom
%  Année universitaire 2025/2026
%  Plateforme intelligente de diagnostic et triage de logs
%  de bases de données hétérogènes par RAG Multi-Moteurs
% ============================================================
\documentclass[12pt,a4paper]{report}

\usepackage[french]{babel}
\usepackage[utf8]{inputenc}
\usepackage[T1]{fontenc}
\usepackage[a4paper,left=2.5cm,right=2.5cm,top=3cm,bottom=2.5cm,headheight=60pt]{geometry}
\usepackage{graphicx}
\usepackage{fancyhdr}
\usepackage{titlesec}
\usepackage[hidelinks,colorlinks=true,linkcolor=blue,citecolor=blue,urlcolor=blue]{hyperref}
\usepackage{xcolor}
\usepackage{enumitem}
\usepackage{caption}
\usepackage{subcaption}
\usepackage{float}
\usepackage{setspace}
\usepackage{tikz}
\usetikzlibrary{shapes.geometric,arrows.meta,positioning,fit,backgrounds}
\usepackage{tabularx}
\usepackage{booktabs}
\usepackage{listings}
\usepackage{lastpage}
\usepackage{amsmath}
\usepackage{microtype}
\usepackage{array}

% ---------- Couleurs ----------
\definecolor{ttblue}{RGB}{31,73,125}
\definecolor{codegreen}{rgb}{0.0,0.5,0.0}
\definecolor{codegray}{rgb}{0.5,0.5,0.5}
\definecolor{backcolor}{rgb}{0.97,0.97,0.97}

% ---------- Style des titres ----------
\titleformat{\chapter}[display]
  {\normalfont\bfseries\Large\color{ttblue}}
  {\Large\chaptertitlename\ \thechapter\ :}{6pt}{\Large}
\titlespacing*{\chapter}{0pt}{0pt}{30pt}
\titleformat{\section}
  {\normalfont\bfseries\large\color{ttblue}}{\thesection.}{0.6em}{}
\titleformat{\subsection}
  {\normalfont\bfseries\normalsize}{\thesubsection.}{0.6em}{}
\titleformat{\subsubsection}
  {\normalfont\bfseries\small}{\thesubsubsection.}{0.6em}{}

% ---------- En-tete / pied de page ----------
\pagestyle{fancy}
\fancyhf{}
\renewcommand{\headrulewidth}{0.6pt}
\fancyhead[L]{\small\textit{Rapport de Stage --- Tunisie Telecom}}
\fancyhead[R]{\small\textit{2025/2026}}
\fancyfoot[C]{\thepage\ / \pageref{LastPage}}
\fancypagestyle{plain}{
  \fancyhf{}
  \fancyhead[L]{\small\textit{Rapport de Stage --- Tunisie Telecom}}
  \fancyhead[R]{\small\textit{2025/2026}}
  \fancyfoot[C]{\thepage\ / \pageref{LastPage}}
  \renewcommand{\headrulewidth}{0.6pt}
}

% ---------- Listings ----------
\lstset{
  backgroundcolor=\color{backcolor},
  basicstyle=\ttfamily\small,
  keywordstyle=\color{ttblue}\bfseries,
  commentstyle=\color{codegreen}\itshape,
  frame=single,breaklines=true,columns=fullflexible,
  showstringspaces=false,
  numbers=left,numberstyle=\tiny\color{codegray},numbersep=5pt
}

% ---------- TikZ styles ----------
\tikzstyle{block}=[rectangle,rounded corners=3pt,minimum width=3.2cm,
  minimum height=0.85cm,text centered,draw=ttblue,fill=blue!5,
  font=\small,text width=3cm]
\tikzstyle{blockW}=[rectangle,rounded corners=3pt,minimum width=4.4cm,
  minimum height=0.85cm,text centered,draw=ttblue,fill=blue!5,
  font=\small,text width=4.2cm]
\tikzstyle{store}=[cylinder,shape border rotate=90,draw=gray!60,fill=gray!10,
  font=\small,minimum height=1.1cm,minimum width=2.8cm,text centered,text width=2.6cm]
\tikzstyle{arrow}=[-{Stealth[length=5pt]},thick,color=ttblue!80]
\tikzstyle{dasharrow}=[-{Stealth[length=5pt]},dashed,thick,color=gray!55]

% ============================================================
\begin{document}

% ============================================================
% PAGE DE GARDE
% ============================================================
\begin{titlepage}
  \begin{minipage}{0.22\textwidth}
    \includegraphics[width=\linewidth]{enit-logo.jpg}
  \end{minipage}
  \hfill
  \begin{minipage}{0.22\textwidth}
    \raggedleft
    \includegraphics[width=\linewidth]{images.jpg}
  \end{minipage}

\thispagestyle{empty}
\begin{center}
\vspace*{0.8cm}
{\large\bfseries\color{ttblue} RAPPORT DE STAGE DE FIN D'\'ETUDES}\\[0.5cm]
\rule{\textwidth}{1.5pt}\\[0.4cm]
{\LARGE\bfseries Plateforme intelligente de diagnostic et triage\\[0.3cm]
de logs de bases de donn\'ees h\'et\'erog\`enes\\[0.3cm]
par G\'en\'eration Augment\'ee par R\'ecup\'eration (RAG) Multi-Moteurs}\\[0.4cm]
\rule{\textwidth}{1.5pt}
\vspace{.6cm}
\begin{tabular}{ll}
\vspace{.2cm}
\textbf{\'Elabor\'e par :}      & Rayen Thabet \\[0.3cm]
\textbf{Classe :}               & 1AMIndS \\[0.3cm]
\textbf{Encadr\'e par :}        & Monsieur Ben Mohamed Ahmed \\[0.15cm]
                                & \small Chef de la Subdivision Administration \\[0.15cm]
                                & \small Syst\`emes et Bases de Donn\'ees \\[0.3cm]
\textbf{Soci\'et\'e d'accueil :}& Tunisie Telecom --- Centre technique IT de la Kasbah \\[0.3cm]
\textbf{D\'epartement :}        & Direction Centrale des Syst\`emes d'Information (DCSI) \\[0.3cm]
\textbf{P\'eriode :}            & Juin -- Ao\^ut 2026 \\[0.3cm]
\textbf{Ann\'ee universitaire :}& 2025/2026 \\
\end{tabular}
\vfill
{\color{ttblue}\rule{0.4\textwidth}{0.8pt}}
\end{center}
\end{titlepage}

% ============================================================
% REMERCIEMENTS
% ============================================================
\chapter*{Remerciements}
\addcontentsline{toc}{chapter}{Remerciements}

Je tiens tout d'abord \`a adresser mes sinc\`eres remerciements \`a l'ensemble du
personnel de Tunisie Telecom, et plus particuli\`erement \`a celui du centre
technique IT de la Kasbah, pour l'accueil chaleureux qui m'a \'et\'e r\'eserv\'e
durant toute la dur\'ee de mon stage.

Mes remerciements les plus sinc\`eres vont \`a mon encadreur, Monsieur Ben Mohamed
Ahmed, Chef de la Subdivision Administration Syst\`emes et Bases de Donn\'ees, qui
m'a accord\'e sa confiance d\`es le premier jour. Sa disponibilit\'e constante, la
pertinence de ses conseils et son expertise technique m'ont \'enorm\'ement apport\'e
et ont grandement contribu\'e \`a la r\'eussite de ce stage. Les \'echanges que nous
avons eus sur les besoins concrets des \'equipes DBA ont directement orient\'e
plusieurs choix techniques du projet.

Je souhaite \'egalement exprimer ma gratitude \`a Monsieur Mustapha Charfeddine
Akkez, dont les conseils avis\'es et le soutien de tous les instants m'ont
accompagn\'e tout au long de cette exp\'erience.

Plus g\'en\'eralement, je remercie chaleureusement toutes les personnes qui, de
pr\`es ou de loin, m'ont fait profiter de leur savoir-faire et de leur exp\'erience
durant ce stage. Travailler dans un environnement aussi exigeant que celui de la
DCSI de Tunisie Telecom a \'et\'e une exp\'erience formatrice et enrichissante,
tant sur le plan technique qu'humain.

% ============================================================
% TABLE DES MATIERES
% ============================================================
\tableofcontents
\listoffigures
\listoftables

% ============================================================
% INTRODUCTION GENERALE
% ============================================================
\chapter*{Introduction g\'en\'erale}
\addcontentsline{toc}{chapter}{Introduction g\'en\'erale}

Les bases de donn\'ees relationnelles constituent l'\'epine dorsale des syst\`emes
informatiques des grandes entreprises. Dans des environnements critiques tels que
ceux de Tunisie Telecom, o\`u Oracle Database, PostgreSQL et MySQL sont exploit\'es
simultan\'ement pour supporter des milliers de transactions par seconde, la
surveillance et le diagnostic rapide des incidents constituent une priorit\'e absolue.

Au c\oe{}ur de cette surveillance se trouve l'analyse des fichiers journaux
(\emph{logs}) produits en continu par chaque instance de base de donn\'ees. Ces
fichiers peuvent atteindre plusieurs dizaines de milliers de lignes par jour,
m\'elant des messages purement informationnels \`a des codes d'erreur critiques
n\'ecessitant une intervention imm\'ediate. Le processus actuel, qui repose sur la
lecture manuelle et l'expertise de l'administrateur de bases de donn\'ees (DBA),
est long, sujet aux omissions et difficile \`a tenir lorsque plusieurs moteurs de
bases de donn\'ees sont administr\'es simultan\'ement.

C'est dans ce contexte que s'inscrit ce stage de fin d'\'etudes, effectu\'e au sein
de la Subdivision Administration Syst\`emes et Bases de Donn\'ees de la Direction
Centrale des Syst\`emes d'Information (DCSI) de Tunisie Telecom, du mois de juin
au mois d'ao\^ut 2026. L'objectif central \'etait de concevoir et d\'evelopper une
plateforme intelligente capable d'analyser automatiquement des fichiers de logs de
trois moteurs de bases de donn\'ees h\'et\'erog\`enes --- Oracle, PostgreSQL et
MySQL ---, d'identifier les incidents critiques, d'en expliquer la signification et
de proposer des actions correctives, le tout sans d\'ependre d'une API de mod\`ele
de langage commerciale externe.

Pour atteindre cet objectif, nous avons adopt\'e l'approche de la
\emph{G\'en\'eration Augment\'ee par R\'ecup\'eration} (RAG,
\emph{Retrieval-Augmented Generation}), qui combine une base de connaissances
structur\'ee \`a un mod\`ele de g\'en\'eration de texte entra\^in\'e sur mesure.

Ce rapport est structur\'e en cinq chapitres. Le \textbf{Chapitre~1} pr\'esente
l'organisme d'accueil. Le \textbf{Chapitre~2} \'etablit la probl\'ematique et
l'\'etat de l'art. Le \textbf{Chapitre~3} expose les d\'ecisions de conception.
Le \textbf{Chapitre~4} d\'etaille la r\'ealisation. Le \textbf{Chapitre~5}
pr\'esente les r\'esultats et une discussion honn\^ete des limites.

% ============================================================
% CHAPITRE 1
% ============================================================
\chapter{Pr\'esentation de l'organisme d'accueil}

\section{Pr\'esentation g\'en\'erale de Tunisie Telecom}

\subsection{Historique et missions}

L'Office National des T\'el\'ecommunications, connu commercialement sous le nom de
\textbf{Tunisie Telecom}, a \'et\'e cr\'e\'e par la loi n\textdegree36 du 17 avril
1995 et est entr\'e en activit\'e le 1\textsuperscript{er} janvier 1996. Fin 2002,
Tunisie Telecom se transforme en soci\'et\'e anonyme de droit public. En mars 1998,
elle inaugure son premier r\'eseau GSM~; en 2009, elle lance le c\^able sous-marin
Hannibal, premier c\^able 100\,\% tunisien~; en 2011, la 3G~; en 2016, la 4G.

Les missions de Tunisie Telecom couvrent l'ensemble de la cha\^ine de valeur des
t\'el\'ecommunications~:
\begin{itemize}[leftmargin=*]
  \item L'installation, le d\'eveloppement, l'entretien et l'exploitation des
        r\'eseaux publics de t\'el\'ecommunications fixes et mobiles~;
  \item La promotion de services \`a valeur ajout\'ee~: internet haut d\'ebit,
        IPTV, solutions cloud pour entreprises~;
  \item La fourniture de solutions d'interconnexion internationale~;
  \item L'h\'ebergement d'infrastructures informatiques critiques (data centers
        certifi\'es Tier~III).
\end{itemize}

\subsection{Chiffres cl\'es et positionnement sur le march\'e tunisien}

\begin{table}[H]
  \centering
  \caption{Chiffres cl\'es de Tunisie Telecom (2025)}
  \label{tab:chiffres_cles}
  \begin{tabularx}{\textwidth}{lX}
    \toprule
    \textbf{Indicateur} & \textbf{Valeur} \\
    \midrule
    Effectif              & Plus de 6\,400 agents \\
    Directions r\'egionales & 24 \\
    Espaces TT            & 140 \\
    Points de vente priv\'es & Plus de 13\,000 \\
    Data centers          & 3 (normes Tier III) \\
    Filiales              & 7 (Sotetel, Topnet, ATI, etc.) \\
    \bottomrule
  \end{tabularx}
\end{table}

\section{Organigramme de Tunisie Telecom}

\subsection{Structure organisationnelle g\'en\'erale}

La figure~\ref{fig:organigramme} pr\'esente la structure organisationnelle de haut
niveau de Tunisie Telecom, avec les directions fonctionnelles rattach\'ees au PDG
et les directions op\'erationnelles sous l'autorit\'e du DGA.

\begin{figure}[H]
  \centering
  \includegraphics[width=0.95\textwidth]{organigramme_tt.png}
  \caption{Organigramme g\'en\'eral de Tunisie Telecom}
  \label{fig:organigramme}
\end{figure}

\subsection{Positionnement de la Direction Centrale des Syst\`emes d'Information}

La \textbf{Direction Centrale des Syst\`emes d'Information (DCSI)} est plac\'ee
directement sous l'autorit\'e du DGA. Elle constitue le bras technique de l'entreprise
pour tout ce qui concerne l'infrastructure informatique, les r\'eseaux internes et
l'administration des bases de donn\'ees.

\section{La Direction Centrale des Syst\`emes d'Information (DCSI)}

\subsection{R\^ole et responsabilit\'es de la DCSI}

La DCSI couvre quatre grands domaines~:
\begin{itemize}[leftmargin=*]
  \item Infrastructure r\'eseau interne (LAN, WAN, intranet)~;
  \item S\'ecurit\'e informatique (firewall, contr\^ole d'acc\`es)~;
  \item Administration des syst\`emes (serveurs Red Hat, Ubuntu, Windows Server)~;
  \item Administration des bases de donn\'ees (Oracle, PostgreSQL, MySQL).
\end{itemize}

\subsection{Le centre technique IT --- Kasbah}

Le centre technique IT de la Kasbah est le site physique o\`u sont h\'eberg\'ees
les principales infrastructures informatiques de Tunisie Telecom. Il abrite les
data centers certifi\'es Tier~III et la Subdivision Administration Syst\`emes et
Bases de Donn\'ees au sein de laquelle j'ai effectu\'e mon stage.

\section{Le service d'accueil et le contexte du stage}

\subsection{La Subdivision Administration Syst\`emes et Bases de Donn\'ees}

Dirig\'ee par Monsieur Ben Mohamed Ahmed (encadreur du stage), cette subdivision
assure~:
\begin{itemize}[leftmargin=*]
  \item La supervision proactive des instances Oracle, PostgreSQL et MySQL~;
  \item L'analyse des journaux d'erreurs~;
  \item La gestion des sauvegardes et restaurations (RMAN, pg\_dump)~;
  \item L'application des correctifs de s\'ecurit\'e.
\end{itemize}

\subsection{Missions quotidiennes de l'\'equipe DBA}

En production, les serveurs Oracle g\'en\`erent des \emph{alert logs} pouvant
atteindre plusieurs milliers de lignes par jour, m\'elant \'ev\'enements b\'enins
et incidents critiques. La lecture de ces fichiers demande une expertise approfondie
et mobilise un temps consid\'erable.

\subsection{Probl\'ematique identifi\'ee \`a l'origine du stage}

Plusieurs observations concordantes ont confirm\'e que \textbf{l'analyse manuelle
des fichiers journaux est un goulot d'\'etranglement op\'erationnel}~:
\begin{itemize}[leftmargin=*]
  \item Trois moteurs (Oracle, PostgreSQL, MySQL) coexistent avec des formats
        de logs et des conventions de codes d'erreur radicalement diff\'erents~;
  \item Un m\^eme code peut appara\^itre des centaines de fois dans un seul fichier~;
  \item Aucun outil interne n'existait pour automatiser ce processus.
\end{itemize}

% ============================================================
% CHAPITRE 2
% ============================================================
\chapter{Contexte et \'etude de l'existant}

\section{Probl\'ematique~: la charge de diagnostic manuel des logs}

\subsection{Le processus actuel de lecture et d'analyse des logs}

Le processus actuel des DBA se d\'eroule en \'etapes manuelles~: ouverture du
fichier via SSH, recherche par \texttt{grep}, consultation de la documentation
pour interpr\'eter les codes, r\'edaction d'un rapport d'incident. Chaque \'etape
est chronophage et non automatis\'ee.

\subsection{Volume et complexit\'e des fichiers de logs en production}

Un serveur Oracle en production g\'en\`ere entre 5\,000 et 50\,000 lignes par
journ\'ee. Dans un environnement multi-instances, le volume total peut d\'epasser
100\,000 lignes quotidiennes.

\subsection{Impact sur le temps des administrateurs de bases de donn\'ees}

L'analyse des logs absorbe plusieurs heures par semaine et par administrateur ---
temps soustrait \`a des t\^aches \`a plus haute valeur ajout\'ee.

\subsection{Absence d'un jeu de donn\'ees r\'eel exploitable (contrainte de confidentialit\'e)}

Les fichiers journaux r\'eels de Tunisie Telecom contiennent des informations
sensibles sur l'architecture interne. La politique de s\'ecurit\'e de la DCSI
interdit formellement leur utilisation en dehors de l'entreprise. Cette contrainte
a conduit \`a une d\'ecision fondamentale~: \textbf{la totalit\'e du corpus de
d\'eveloppement a \'et\'e g\'en\'er\'ee synth\'etiquement}, reproduisant
fid\`element les formats r\'eels sans utiliser de donn\'ees de production.

\section{Les syst\`emes cibl\'es et les limites de l'analyse manuelle actuelle}

\subsection{Oracle Database --- format des logs et codes d'erreur}

Les journaux Oracle (alert log) suivent un format texte structur\'e, horodat\'e
ISO~8601. Les codes respectent une convention de pr\'efixage~:

\begin{table}[H]
  \centering
  \caption{Principaux pr\'efixes d'erreurs Oracle}
  \label{tab:prefixes_oracle}
  \begin{tabularx}{\textwidth}{lX}
    \toprule
    \textbf{Pr\'efixe} & \textbf{Sous-syst\`eme} \\
    \midrule
    \texttt{ORA-}  & Erreurs g\'en\'erales du noyau Oracle \\
    \texttt{TNS-}  & Couche r\'eseau Oracle Net \\
    \texttt{RMAN-} & Recovery Manager \\
    \texttt{PLS-}  & Moteur PL/SQL \\
    \texttt{CRS-}  & Oracle Clusterware (RAC) \\
    \bottomrule
  \end{tabularx}
\end{table}

\begin{lstlisting}[language={},caption={Extrait type d'un alert log Oracle},label=lst:alertlog_oracle]
2026-07-01T08:13:42.312041+01:00
ORA-01555 encountered when current SCN is 12304568
ORA-01555: snapshot too old: rollback segment number 4 with name "RBS3"
\end{lstlisting}

\subsection{PostgreSQL --- format des logs et codes SQLSTATE}

PostgreSQL utilise des codes SQLSTATE \`a cinq caract\`eres (ex.~\texttt{42P01})
et des niveaux de s\'ev\'erit\'e explicites~: \texttt{ERROR}, \texttt{FATAL},
\texttt{PANIC}, \texttt{WARNING}.

\begin{lstlisting}[language={},caption={Extrait type d'un log PostgreSQL},label=lst:alertlog_pg]
2026-07-02 10:14:32.112 CET [12345] ERROR:  relation "users_tmp" does not exist
2026-07-02 10:14:32.113 CET [12345] STATEMENT: SELECT * FROM users_tmp;
\end{lstlisting}

\subsection{MySQL --- format structur\'e et codes d'erreur}

Depuis MySQL~8, le journal adopte un format structur\'e avec des balises
\texttt{[MY-XXXXXX]} et des sous-syst\`emes \texttt{[Server]}, \texttt{[InnoDB]}.

\begin{lstlisting}[language={},caption={Extrait type d'un log MySQL~8},label=lst:alertlog_mysql]
2026-07-03T14:22:11.567890Z 0 [ERROR] [MY-010065] [Server] Could not create
  unix socket lock file /var/run/mysqld/mysqld.sock.lock.
\end{lstlisting}

\subsection{\'Etude comparative des formats et limites de l'analyse manuelle}

\begin{table}[H]
  \centering
  \caption{Comparaison des formats de logs des trois moteurs}
  \label{tab:compare_formats}
  \begin{tabularx}{\textwidth}{lXXX}
    \toprule
    \textbf{Crit\`ere} & \textbf{Oracle} & \textbf{PostgreSQL} & \textbf{MySQL~8+} \\
    \midrule
    Format code        & Pr\'efixe-num\'ero & SQLSTATE 5~chars & {[MY-XXXXXX]} \\
    S\'ev\'erit\'e     & Implicite & Explicite & Explicite \\
    Horodatage         & ISO~8601 & Configurable & ISO~8601 UTC \\
    Codes info.        & Fr\'equents & Rares & Fr\'equents \\
    \bottomrule
  \end{tabularx}
\end{table}

\section{Limites des approches classiques}

\subsection{L'analyse manuelle ligne par ligne}

Impraticable au-del\`a de quelques centaines de lignes. Requiert une expertise
que les membres moins exp\'eriment\'es de l'\'equipe ne poss\`edent pas.

\subsection{Les scripts de recherche par mots-cl\'es (grep, regex simples)}

Utiles pour l'extraction, mais ne fournissent aucune explication ni recommandation.
Ne distinguent pas les erreurs r\'eelles des messages informationnels.

\subsection{Les outils de log management (ELK, Splunk) et leurs limites}

ELK et Splunk offrent de bonnes capacit\'es de visualisation, mais ne g\'en\`erent
pas d'explication textuelle du sens d'un code d'erreur ni de recommandation d'action
corrective. Ils n\'ecessitent \'egalement une infrastructure d\'edi\'ee.

\section{\'Etat de l'art~: Intelligence Artificielle appliqu\'ee \`a l'analyse de logs}

\subsection{D\'etection d'anomalies dans les logs (LogBERT, NeuralLog)}

LogBERT et NeuralLog, \'etudi\'es en d\'etail dans les benchmarks de r\'ef\'erence~\cite{le2022log,zhu2019log},
repr\'esentent les travaux majeurs d'application des architectures Transformer aux journaux syst\`eme.
Ces approches sont toutefois orient\'ees \textbf{d\'etection binaire} (s\'equence normale vs.
anormale) et ne produisent pas d'explication textuelle ni de recommandation d'action --- ce qui les rend
insuffisantes pour le cas d'usage cible des DBA.

\subsection{Le paradigme RAG (Retrieval-Augmented Generation)}

La G\'en\'eration Augment\'ee par R\'ecup\'eration (RAG), introduite par Lewis et al.~\cite{lewis2020rag},
associe une \'etape de recherche documentaire \`a une \'etape de g\'en\'eration de
texte. Elle ancre les r\'eponses sur des documents pr\'ealablement s\'electionn\'es,
r\'eduisant les hallucinations et am\'eliorant la tra\c cabilit\'e des diagnostics.

\subsection{Les mod\`eles de langage pr\'e-entra\^in\'es et le fine-tuning}

Le fine-tuning d'un mod\`ele comme FLAN-T5~\cite{flanT5} sur un corpus de domaine
permet d'adapter ses capacit\'es g\'en\'eratives \`a une t\^ache pr\'ecise. Avec
un corpus d'entra\^inement relativement modeste, on obtient des performances
satisfaisantes sur une t\^ache de g\'en\'eration de texte structur\'e.

\subsection{Comparatif des architectures~: encodeur-d\'ecodeur vs d\'ecodeur seul}

\begin{itemize}[leftmargin=*]
  \item \textbf{Encodeur-d\'ecodeur} (T5, BART)~: l'encodeur traite l'entr\'ee
        et le d\'ecodeur g\'en\`ere la sortie structur\'ee. Adapт\'e aux t\^aches
        de g\'en\'eration conditionnée~;
  \item \textbf{D\'ecodeur seul} (GPT, LLaMA)~: tr\`es performants en g\'en\'eration
        libre, mais plus difficiles \`a contraindre \`a un format de sortie rigide.
\end{itemize}

\section{Choix strat\'egiques et justification}

\subsection{Pourquoi \'eviter le recours \`a une API LLM payante}

Les API commerciales pr\'esentent des inconv\'enients r\'edhibitoires~: les logs
ne doivent pas transiter par des serveurs tiers~; le co\^ut est r\'ecurrent~; la
disponibilit\'e d\'epend d'une infrastructure externe.

\subsection{Le choix d'un mod\`ele entra\^in\'e localement}

FLAN-T5 base (250M param\`etres) est suffisamment compact pour fonctionner sur CPU
cloud sans GPU, compatible avec les contraintes de d\'eploiement de Railway.

\subsection{Le choix de l'architecture RAG plut\^ot qu'un mod\`ele purement g\'en\'eratif}

Sans RAG, le mod\`ele risque de produire des explications plausibles mais incorrectes
pour des codes rares. Avec RAG, si le code est absent de la base, le syst\`eme le
signale explicitement plut\^ot que d'inventer.

\subsection{Synth\`ese des choix techniques retenus}

\begin{table}[H]
  \centering
  \caption{Synth\`ese des choix techniques retenus}
  \label{tab:choix_techniques}
  \begin{tabularx}{\textwidth}{lX}
    \toprule
    \textbf{Composant} & \textbf{Technologie retenue et justification} \\
    \midrule
    G\'en\'eration   & FLAN-T5 base fine-tun\'e~: l\'eger, encodeur-d\'ecodeur, Apache~2.0 \\
    R\'ecup\'eration & BM25~: sans mod\`ele externe, adapt\'e aux textes courts et techniques \\
    Classification   & TF-IDF + R\'egression Logistique~: interpr\'etable, rapide \\
    Backend web      & FastAPI~: performant, typage fort, documentation automatique \\
    D\'eploiement    & Docker + Railway + Hugging Face Hub \\
    \bottomrule
  \end{tabularx}
\end{table}

% ============================================================
% CHAPITRE 3
% ============================================================
\chapter{Conception de la solution}

\section{Architecture g\'en\'erale du pipeline}

\subsection{Vue d'ensemble}

La solution est un pipeline de traitement s\'equentiel \`a six \'etapes~:
\textbf{d\'etection du moteur} \textrightarrow{} \textbf{extraction et d\'eduplication}
\textrightarrow{} \textbf{classification} \textrightarrow{} \textbf{r\'ecup\'eration}
\textrightarrow{} \textbf{g\'en\'eration} \textrightarrow{} \textbf{rapport}.
Chaque \'etape est d\'ecouple\'ee et testable ind\'ependamment.

\subsection{Sch\'ema d'architecture d\'etaill\'e}

\begin{figure}[H]
\centering
\begin{tikzpicture}[node distance=0.55cm and 2.2cm]
  \node[blockW] (log)     {\textbf{Fichier journal}\\{\small Oracle / PG / MySQL}};
  \node[blockW, below=of log]    (detect)  {\textbf{D\'etection moteur}\\{\small engine\_detection.py}};
  \node[blockW, below=of detect] (parser)  {\textbf{Parseur multi-moteurs}\\{\small Oracle, PG, MySQL}};
  \node[blockW, below=of parser] (dedup)   {\textbf{D\'eduplication}\\{\small codes uniques + comptages}};
  \node[blockW, below=of dedup]  (classif) {\textbf{Classifieur}\\{\small TF-IDF + LR}};
  \node[store, right=2.4cm of parser]  (kb) {\textbf{Base de}\\{\textbf{connaissances}}\\{\small 27\,622 entr\'ees}};
  \node[blockW, below=1.6cm of kb] (ret) {\textbf{R\'ecup\'erateur BM25}\\{\small exact (score=999)\\ puis lexical}};
  \node[blockW, below=of ret]    (gen)      {\textbf{G\'en\'eration FLAN-T5}\\{\small fine-tun\'e multi-moteurs}};
  \node[blockW, below=of gen]    (report)   {\textbf{Rapport}\\{\small JSON + Markdown + Web}};
  \draw[arrow] (log) -- (detect);
  \draw[arrow] (detect) -- (parser);
  \draw[arrow] (parser) -- (dedup);
  \draw[arrow] (dedup) -- (classif);
  \draw[arrow] (classif.east) -| (ret.west);
  \draw[arrow] (kb) -- (ret);
  \draw[arrow] (ret) -- (gen);
  \draw[arrow] (gen) -- (report);
\end{tikzpicture}
\caption{Architecture d\'etaill\'ee du pipeline RAG multi-moteurs}
\label{fig:pipeline-global}
\end{figure}

\subsection{Flux de donn\'ees entre les composants}

Le flux suit un chemin lin\'eaire~: le fichier est soumis au d\'etecteur de moteur,
puis au parseur appropri\'e, qui extrait les codes avec leur contexte. Apr\`es
d\'eduplication, chaque code unique est soumis au classifieur, puis au r\'ecup\'erateur
qui interroge la base de connaissances, et enfin au g\'en\'erateur qui produit
l'explication structur\'ee.

\section{Conception de la base de connaissances}

\subsection{Structure et sch\'ema des entr\'ees}

\begin{lstlisting}[language={},caption={Sch\'ema d'une entr\'ee de la base de connaissances},label=lst:kb_schema]
{
  "code":     "ORA-01555",
  "message":  "snapshot too old",
  "category": "Transaction",
  "cause":    "Segment de rollback recycle avant fin de lecture longue.",
  "solution": "Augmenter UNDO_RETENTION ou la taille des segments UNDO.",
  "keywords": ["snapshot", "rollback", "undo", "scn"],
  "severity": "high",
  "engine":   "oracle"
}
\end{lstlisting}

\subsection{Strat\'egie de construction pour Oracle (documentation officielle)}

Construction en deux phases~: 58 entr\'ees manuelles (codes les plus fr\'equents),
puis 27\,224 entr\'ees suppl\'ementaires par web scraping de la documentation
Oracle~\cite{oracledoc} via BeautifulSoup~\cite{bs4}.

\subsection{Strat\'egie de construction pour PostgreSQL (fichier source SQLSTATE)}

Le fichier \texttt{errcodes.txt} du code source officiel PostgreSQL a \'et\'e
pars\'e~: 261 codes SQLSTATE standards avec leur libel\'e officiel.

\subsection{Strat\'egie de construction pour MySQL (codes cur\'es manuellement)}

79 codes MySQL les plus fr\'equents, saisis manuellement depuis la documentation
officielle MySQL~8. La faible couverture est une limite reconnue du projet.

\subsection{Fusion en une base de connaissances multi-moteurs unifi\'ee}

Les trois bases ont \'et\'e fusionn\'ees en un fichier JSONL unique
(\texttt{combined\_errors\_kb.jsonl}). Le champ \texttt{engine} permet le filtrage
par moteur pour \'eviter la contamination inter-moteurs. Total~: \textbf{27\,622 entr\'ees}
(9,78~Mo).

\section{Strat\'egie de g\'en\'eration de donn\'ees synth\'etiques}

\subsection{Justification~: absence de logs r\'eels exploitables}

La totalit\'e du corpus d'entra\^inement est synth\'etique (cf.\ section~2.1.4).

\subsection{G\'en\'eration de logs synth\'etiques r\'ealistes par moteur}

Le script \texttt{generate\_finetune\_logs\_v2.py} produit des fichiers journaux
reproduisant fid\`element la structure des vrais logs~: horodatages ISO~8601,
lignes de bruit intercal\'ees (archivage, checkpoint), injection al\'eatoire de
codes d'erreur.

\subsection{Diversit\'e des occurrences et gestion des valeurs de substitution}

Les placeholders (\texttt{\%s}, \texttt{string}) sont r\'esolus avant g\'en\'eration
en injectant des valeurs r\'ealistes (noms de tablespace, num\'eros de segment, etc.)
pour produire des messages d'erreur plausibles.

\section{Conception du module de r\'ecup\'eration}

\subsection{Choix de BM25 plut\^ot que les embeddings denses}

BM25~\cite{bm25orig} a \'et\'e retenu car~: (1)~sans mod\`ele externe, fonctionnement
hors ligne~; (2)~les textes Oracle/PG/MySQL sont courts et riches en termes techniques
o\`u le lexical est tr\`es efficace~; (3)~latence de l'ordre de quelques dizaines
de millisecondes sur 27\,000 entr\'ees.

\subsection{Recherche par correspondance exacte de code}

Avant toute recherche lexicale, une correspondance exacte sur le code normalis\'e
est tent\'ee. Si trouv\'ee, l'entr\'ee est retourn\'ee avec un score sentinelle de~999.
Cette correspondance court-circuite le mod\`ele T5.

\subsection{Filtrage par moteur pour \'eviter la contamination inter-moteurs}

Les r\'esultats BM25 sont filtr\'es pour ne retenir que les entr\'ees du moteur
d\'etect\'e (oracle, postgres, mysql), \'evitant des correspondances s\'emantiquement
incorrectes.

\subsection{Gestion des codes non identifi\'es (pseudo-codes)}

Les messages sans code normalis\'e (ex.~\texttt{ERROR: could not connect}) re\c coivent
un pseudo-code (\texttt{PG-ERROR-xxx}). La correspondance exacte est d\'esactiv\'ee
pour ces pseudo-codes.

\section{Conception du module de classification}

\subsection{Probl\'ematique~: distinguer erreur r\'eelle et message informationnel}

Certains codes (ex.~ORA-16111) sont de simples messages d'\'etat. Les envoyer au
g\'en\'erateur serait un gaspillage de ressources et induirait les op\'erateurs en erreur.

\subsection{Choix d'un mod\`ele simple~: TF-IDF + r\'egression logistique}

Le classifieur combine deux matrices TF-IDF (unigrammes/bigrammes de mots et
n-grammes de caract\`eres) concat\'en\'ees avant une r\'egression logistique avec
pond\'eration de classes.

\subsection{Strat\'egie d'\'etiquetage automatique \`a partir de la base de connaissances}

Les entr\'ees KB mentionnant explicitement leur caract\`ere informatif sont
\'etiquet\'ees \texttt{INFORMATIONAL}. Cette heuristique a identifi\'e 878 codes
Oracle informationnels sur 27\,282.

\section{Choix et justification du mod\`ele g\'en\'eratif}

\subsection{Pourquoi FLAN-T5 plut\^ot qu'un mod\`ele d\'ecodeur seul}

FLAN-T5~\cite{flanT5} est naturellement contraint \`a condenser l'information de
l'entr\'ee plut\^ot qu'\`a la prolonger librement, ce qui limite les hallucinations.

\subsection{Le principe de l'ancrage (\emph{grounding})}

Le prompt envoy\'e au mod\`ele T5 contient syst\'ematiquement les entr\'ees KB
r\'ecup\'er\'ees pour le code analys\'e, garantissant que les explications sont
ancr\'ees sur des donn\'ees v\'erifi\'ees.

\subsection{Conception du format d'entra\^inement}

\begin{lstlisting}[language={},caption={Format d'entr\'ee du mod\`ele T5},label=lst:t5_format]
Analyze this Oracle Database error log entry and explain the error.

ORA-01555: snapshot too old: rollback segment number 4 with name "RBS3"

RETRIEVED KNOWLEDGE:
- [ORA-01555] snapshot too old
  cause: Segment de rollback recycle avant fin de lecture longue.
  solution: Augmenter UNDO_RETENTION ou la taille des segments UNDO.
  (score: 999.0)
\end{lstlisting}

% ============================================================
% CHAPITRE 4
% ============================================================
\chapter{R\'ealisation et impl\'ementation}

\section{Extraction et analyse des logs (parseurs)}

\subsection{Impl\'ementation du parseur Oracle}

Le parseur Oracle lit le fichier journal ligne par ligne et applique des expressions
r\'eguli\`eres couvrant \textbf{81 pr\'efixes} reconnus (\'etendu de 13 \`a 81
apr\`es tests). Pour chaque occurrence, il conserve le code normalis\'e, le num\'ero
de ligne, l'horodatage et une fen\^etre de contexte configurable.

\subsection{Impl\'ementation du parseur PostgreSQL}

G\`ere plusieurs formats (texte avec ou sans CSV) et extrait les niveaux de
s\'ev\'erit\'e. Les messages sans SQLSTATE re\c coivent un pseudo-code
(\texttt{PG-FATAL}, \texttt{PG-ERROR}).

\subsection{Impl\'ementation du parseur MySQL et gestion des formats h\'erit\'es}

G\`ere le format structur\'e MySQL~8 (balises \texttt{[MY-XXXXXX]}) et le format
textuel h\'erit\'e des versions ant\'erieures. La d\'etection du format est automatique.

\subsection{Normalisation des codes d'erreur}

La fonction \texttt{normalize\_code()} harmonise les variantes syntaxiques~:
\texttt{ORA-1652} et \texttt{ORA-01652} d\'esignent la m\^eme erreur (z\'ero-padding
sur 5 chiffres pour Oracle).

\section{D\'etection automatique du moteur de base de donn\'ees}

\subsection{Approche par heuristiques et vote pond\'er\'e}

Le module \texttt{src/engine\_detection.py} applique un syst\`eme de vote
pond\'er\'e~: pr\'esence de \texttt{ORA-}/\texttt{TNS-} (Oracle), de niveaux
\texttt{FATAL}/\texttt{PANIC} et codes SQLSTATE (PostgreSQL), de balises
\texttt{[MY-XXXXXX]}/\texttt{[InnoDB]} (MySQL).

\subsection{Calcul d'un score de confiance}

Un score normalis\'e entre 0 et 1 est calcul\'e pour chaque moteur. Le moteur avec
le score le plus \'elev\'e est s\'electionn\'e. Ce score est inclus dans le rapport final.

\subsection{Gestion des cas ambigus et des logs \`a faible signal}

Si le score maximum est inf\'erieur \`a 0,1, un avertissement est \'emis dans le
rapport et le syst\`eme se replie sur Oracle par d\'efaut.

\section{Construction et enrichissement de la base de connaissances}

\subsection{Extraction automatis\'ee depuis la documentation Oracle}

Le script \texttt{build\_knowledge\_base.py} extrait les codes Oracle via
BeautifulSoup~\cite{bs4}~: code normalis\'e, message court, cause, actions recommand\'ees.
Production~: 27\,224 nouvelles entr\'ees venant compl\'eter les 58 initiales.

\subsection{Extraction depuis le fichier source officiel PostgreSQL}

Parsing du fichier \texttt{errcodes.txt}~: 261 codes SQLSTATE standards.

\subsection{Curation manuelle des codes MySQL les plus fr\'equents}

79 codes MySQL saisis manuellement depuis la documentation officielle MySQL~8.

\subsection{Statistiques finales de couverture de la base de connaissances}

\begin{table}[H]
  \centering
  \caption{Statistiques de la base de connaissances unifi\'ee}
  \label{tab:kb_stats}
  \begin{tabular}{lrrr}
    \toprule
    \textbf{Moteur} & \textbf{Entr\'ees} & \textbf{Taille (Mo)} & \textbf{M\'ethode} \\
    \midrule
    Oracle      & 27\,282 & 9,12 & Web scraping + manuel \\
    PostgreSQL  & 261     & 0,11 & Fichier source officiel \\
    MySQL       & 79      & 0,04 & Curation manuelle \\
    \midrule
    \textbf{Total} & \textbf{27\,622} & \textbf{9,78} & --- \\
    \bottomrule
  \end{tabular}
\end{table}

\section{D\'eveloppement du classifieur}

\subsection{G\'en\'eration du jeu de donn\'ees d'entra\^inement \'etiquet\'e}

Chaque occurrence est \'etiquet\'ee \texttt{REAL\_ERROR} ou \texttt{INFORMATIONAL}
depuis la KB. Le d\'es\'equilibre initial (2,9\,\% d'informationnels) est corrig\'e
par surrepr\'esentation lors de la g\'en\'eration, portant la proportion \`a 17,5\,\%.

\subsection{Entra\^inement et comparaison de mod\`eles}

R\'egression logistique (TF-IDF) et SVM lin\'eaire compar\'es~; la r\'egression
logistique est retenue pour ses probabilit\'es calibr\'ees.

\subsection{Calibrage du seuil de d\'ecision}

Seuil calibr\'e pour un rappel de~0,90 sur la classe informationnelle.

\subsection{R\'esultats et limites du classifieur}

\begin{table}[H]
  \centering
  \caption{Performances du classifieur sur le jeu de test}
  \label{tab:classif_perf}
  \begin{tabular}{lccc}
    \toprule
    \textbf{Classe} & \textbf{Pr\'ecision} & \textbf{Rappel} & \textbf{F1} \\
    \midrule
    Informationnelle & 0,863 & 0,901 & \textbf{0,882} \\
    Erreur r\'eelle  & 0,994 & 0,974 & 0,984 \\
    \bottomrule
  \end{tabular}
\end{table}

\section{Fine-tuning du mod\`ele FLAN-T5}

\subsection{G\'en\'eration du corpus de fine-tuning multi-moteurs}

Pour chaque occurrence dans les logs synth\'etiques, un exemple d'entra\^inement
est construit~: le prompt contient l'instruction sp\'ecifique au moteur, le contexte
du log et les entr\'ees KB r\'ecup\'er\'ees~; la cible est la sortie structur\'ee
\`a quatre champs.

\subsection{S\'eparation train / validation / test par code}

La s\'eparation s'effectue \textbf{strictement par code d'erreur} et non par ligne de log.
Cela garantit que le jeu de test contient exclusivement des codes totalement absents du
jeu d'entra\^inement, permettant d'\'evaluer sans biais la capacit\'e de g\'en\'eralisation
du mod\`ele \`a des anomalies inconnues.

Le corpus v2 multi-moteurs int\`egre les trois moteurs de bases de donn\'ees. Le
tableau~\ref{tab:splits} pr\'esente la r\'epartition exacte par moteur~:

\begin{table}[H]
  \centering
  \caption{R\'epartition d\'etaill\'ee du corpus de fine-tuning multi-moteurs (v2)}
  \label{tab:splits}
  \begin{tabularx}{\textwidth}{lrrrr}
    \toprule
    \textbf{Sous-ensemble} & \textbf{Oracle} & \textbf{PostgreSQL} & \textbf{MySQL} & \textbf{Total (Codes uniques)} \\
    \midrule
    Entra\^inement (train)      & 85\,383 & 840 & 200 & \textbf{86\,423} (23\,500) \\
    Validation (val)            & 10\,005 &  78 &  42 & \textbf{10\,125} (\phantom{0}2\,771) \\
    Test (codes non vus)        &  4\,896 &  72 &  24 &  \textbf{4\,992} (\phantom{0}1\,348) \\
    \midrule
    \textbf{Total g\'en\'eral}  & \textbf{100\,284} & \textbf{990} & \textbf{266} & \textbf{101\,540} (27\,619) \\
    \bottomrule
  \end{tabularx}
\end{table}

Comme le met en \'evidence le tableau~\ref{tab:splits}, la proportion d'exemples Oracle
repr\'esente 98,7\,\% du corpus total, refl\'etant la disparit\'e de taille des bases de
connaissances sources. Cette asym\'etrie assum\'ee constitue l'une des limites analys\'ees
en section~\ref{sec:limites}.

\subsection{Environnement et hyperparam\`etres d'entra\^inement}

\begin{itemize}[leftmargin=*]
  \item Mod\`ele de base~: \texttt{google/flan-t5-base} (250M param\`etres)~;
  \item Mode~: \emph{full fine-tuning} (sans LoRA), sur Kaggle (GPU Tesla T4, 16~Go)~;
  \item \'Epoques~: 4~; Batch size~: 16~; Max tokens entr\'ee~: 512, sortie~: 200~;
  \item Optimiseur~: AdamW avec warmup de 500 pas~;
  \item Dur\'ee totale~: environ 6h30.
\end{itemize}

\subsection{Int\'egration de l'ancrage par connaissances r\'ecup\'er\'ees}

Chaque exemple inclut syst\'ematiquement les entr\'ees KB r\'ecup\'er\'ees (jusqu'\`a
3). Cela force le mod\`ele \`a apprendre \`a utiliser les connaissances fournies.

\subsection{D\'etection et correction d'un bug de mesure critique}

L'indicateur d'adh\'erence au format \'etait bloqu\'e \`a z\'ero sur toutes les
\'epoques. Le bug identifi\'e~: les tokeniseurs T5 suppriment les sauts de ligne \`a
l'encodage, rendant le parseur de r\'eponse initial inop\'erant. Correction~:
remplacement du d\'ecoupage ligne par ligne par une regex avec \emph{lookahead}.
Apr\`es correction~: adh\'erence \textbf{> 90\,\%} sur le jeu de validation.

\subsection{Calibrage du niveau de confiance dans les sorties}

Seules les correspondances exactes re\c coivent \texttt{CONFIDENCE: high}. Les
correspondances BM25 re\c coivent \texttt{medium} ou \texttt{low} selon le score.

\section{Int\'egration du mod\`ele dans le pipeline applicatif}

\subsection{Chargement et mise en cache du mod\`ele}

Le mod\`ele est h\'eberg\'e sur Hugging Face Hub
(\texttt{rayenthabet004/tt-multi-engine-t5}). Les poids sont pr\'e-t\'el\'echarg\'es
\`a la construction de l'image Docker (\texttt{ENV HF\_HOME=/app/hf\_cache}).
Un dictionnaire \texttt{\_T5\_MODEL\_CACHE} garantit un chargement unique par instance.

\subsection{Court-circuit d\'eterministe pour les correspondances exactes}

Pour les codes pr\'esents dans la base (score~= 999), la g\'en\'eration T5 est
court-circuit\'ee. L'explication est construite directement depuis la KB.

\subsection{Garde-fous post-g\'en\'eration}

(1)~Filtre de r\'ep\'etition~: d\'etecte et supprime les n-grammes r\'ep\'et\'es.
(2)~Filtre de coh\'erence~: d\'etecte les codes erron\'es cit\'es dans la r\'eponse
(hallucination) et les remplace par des formulations neutres.

\section{D\'eveloppement de l'application web}

\subsection{Architecture backend (FastAPI)}

Principaux endpoints~:
\begin{itemize}[leftmargin=*]
  \item \texttt{POST /api/analyze}~: analyse d'un log en JSON~;
  \item \texttt{POST /api/analyze/upload}~: analyse d'un fichier upload\'e~;
  \item \texttt{POST /api/chat}~: chatbot DBA interactif~;
  \item \texttt{GET /api/health}~: \'etat du service.
\end{itemize}

Optimisation critique~: la base de connaissances et l'index BM25 sont
\textbf{construits une seule fois au d\'emarrage} et r\'eutilis\'es via des
singletons (\texttt{\_KB}, \texttt{\_RETRIEVER}). Sans cette optimisation, chaque
requ\^ete ajoutait 2 \`a 4~secondes de latence.

\subsection{Conception de l'interface utilisateur}

Interface HTML/CSS/JS vanilla aux couleurs Tunisie Telecom (bleu nuit, cyan).
Fonctionnalit\'es~: tableau de bord KPIs, mode clair/sombre, filtres par moteur
et s\'ev\'erit\'e, export Markdown et JSON.

\subsection{Fonctionnalit\'e de chat interactif sur les rapports g\'en\'er\'es}

Assistant conversationnel int\'egr\'e h\'eritant automatiquement du contexte du
dernier rapport. Guardrails stricts rejetant les questions hors du domaine DBA.

\section{D\'eploiement}

\subsection{Conteneurisation avec Docker}

\begin{lstlisting}[language={},caption={Extrait du Dockerfile --- pr\'e-t\'el\'echargement du mod\`ele},label=lst:dockerfile]
ENV HF_HOME=/app/hf_cache
RUN python -c "from transformers import AutoTokenizer, AutoModelForSeq2SeqLM; \
    AutoTokenizer.from_pretrained('rayenthabet004/tt-multi-engine-t5'); \
    AutoModelForSeq2SeqLM.from_pretrained('rayenthabet004/tt-multi-engine-t5', \
    low_cpu_mem_usage=True)"
ENV HF_HUB_OFFLINE=1
\end{lstlisting}

\subsection{H\'ebergement des poids du mod\`ele (Hugging Face Hub)}

Les poids (\textasciitilde{}950~Mo) sont h\'eberg\'es sous l'identifiant
\texttt{rayenthabet004/tt-multi-engine-t5}, int\'egr\'es automatiquement au build
Docker sans alourdir le d\'ep\^ot Git.

\subsection{D\'eploiement sur Railway}

Railway d\'etecte automatiquement le Dockerfile et d\'eploie \`a chaque push sur
\texttt{main}. URL de production~:
\url{https://projet-tt-production.up.railway.app}

\subsection{Probl\`emes rencontr\'es et r\'esolutions}

Trois probl\`emes de d\'eploiement ont \'et\'e identifi\'es et r\'esolus~:

\begin{enumerate}[leftmargin=*]
  \item \textbf{Chargement infini (bug reload=True)}~: le file watcher Uvicorn
        d\'etectait les fichiers de cache HF comme des modifications et red\'emarrait
        le serveur en cours de requ\^ete. \textit{Solution}~: \texttt{reload}
        conditionnel \`a \texttt{ENV=development}~;

  \item \textbf{Reconstruction KB par requ\^ete}~: la base et l'index BM25 \'etaient
        reconstruits \`a chaque appel HTTP (+2 \`a 4~s de latence).
        \textit{Solution}~: mise en cache en singleton au d\'emarrage~;

  \item \textbf{OOM caus\'e par torch.compile()}~: TorchInductor ajoutait 300--500~Mo
        de cache noyaux compil\'es en plus du mod\`ele T5 (1,1~Go), provoquant un
        arr\^et par manque de m\'emoire.
        \textit{Solution}~: abandon de \texttt{torch.compile()}.
\end{enumerate}

% ============================================================
% CHAPITRE 5
% ============================================================
\chapter{Tests, r\'esultats et \'evaluation}

\section{M\'ethodologie de test}

\subsection{Tests unitaires et d'int\'egration}

La suite comprend \textbf{55~tests} r\'epartis en cinq fichiers~:

\begin{table}[H]
  \centering
  \caption{R\'epartition des tests automatis\'es}
  \label{tab:tests}
  \begin{tabular}{lr}
    \toprule
    \textbf{Fichier de test} & \textbf{Tests} \\
    \midrule
    \texttt{test\_engine\_and\_parsers.py} & 29 \\
    \texttt{test\_log\_parser.py}          & 7 \\
    \texttt{test\_multi\_engine\_e2e.py}   & 8 \\
    \texttt{test\_pipeline.py}             & 5 \\
    \texttt{test\_retriever.py}            & 6 \\
    \midrule
    \textbf{Total}                         & \textbf{55} \\
    \bottomrule
  \end{tabular}
\end{table}

Tous les tests sont ex\'ecutables sans connexion internet ni cl\'e d'API.

\subsection{G\'en\'eration de logs de test synth\'etiques cibl\'es}

Trois fichiers \emph{challenging} ont \'et\'e cr\'e\'es manuellement~:
\texttt{oracle\_challenging.log}, \texttt{postgresql\_challenging.log} et
\texttt{mysql\_challenging.log}. Ils contiennent des codes rares, des codes
informationnels et des combinaisons susceptibles de tester les cas limites.

\subsection{Tests sur des sc\'enarios r\'ealistes et volontairement complexes}

Ces fichiers sont utilis\'es pour des tests de non-r\'egression \`a chaque
modification significative du pipeline.

\section{R\'esultats quantitatifs}

\subsection{Pr\'ecision et rappel du classifieur}

F1 de 0,882 sur la classe informationnelle (satisfaisant compte tenu du
d\'es\'equilibre initial). F1 de 0,984 sur la classe erreur r\'eelle.

\subsection{Adh\'erence au format de sortie structur\'e}

L'\'evaluation de l'adh\'erence syntaxique (pr\'esence des quatre champs \texttt{MEANING},
\texttt{LIKELY\_CAUSE}, \texttt{SUGGESTED\_SOLUTION}, \texttt{CONFIDENCE} correctement
d\'elimit\'es) d\'emontre une excellente robustesse du mod\`ele T5 apr\`es fine-tuning~:
le taux d'adh\'erence global atteint \textbf{98,9\,\%} sur l'ensemble du jeu de test
(97,0\,\% pour Oracle, 100\,\% pour PostgreSQL et 100\,\% pour MySQL).

\subsection{Scores BLEU-4 par moteur et analyse de la profondeur des connaissances}

L'\'evaluation quantitative de la qualit\'e des diagnostics g\'en\'er\'es a \'et\'e mesur\'ee
\`a l'aide du score \textbf{BLEU-4} calcul\'e par rapport aux explications de r\'ef\'erence
sur des codes d'erreur strictement non vus lors de l'entra\^inement. Le tableau~\ref{tab:bleu}
pr\'esente les r\'esultats obtenus par moteur~:

\begin{table}[H]
  \centering
  \caption{M\'etriques d'\'evaluation quantitative par moteur (jeu de test, codes non vus)}
  \label{tab:bleu}
  \begin{tabularx}{\textwidth}{lrrr}
    \toprule
    \textbf{Moteur} & \textbf{BLEU-4} & \textbf{Adh\'erence format (\%)} & \textbf{Taille KB (entr\'ees)} \\
    \midrule
    Oracle          & \textbf{0,210}  & 97,0\,\%                         & 27\,282 (compl\`ete) \\
    MySQL           & 0,170           & 100,0\,\%                        & 79 (cur\'ee manuellement) \\
    PostgreSQL      & 0,080           & 100,0\,\%                        & 261 (38 cur\'ees, 223 g\'en\'eriques) \\
    \midrule
    \textbf{Moyenne globale} & \textbf{0,180} & \textbf{98,9\,\%}       & \textbf{27\,622} \\
    \bottomrule
  \end{tabularx}
\end{table}

\subsection{Analyse critique des r\'esultats BLEU-4 et limites des m\'etriques n-grammes}

L'analyse du tableau~\ref{tab:bleu} r\'ev\`ele un enseignement technique fondamental~:
\textbf{PostgreSQL obtient le score BLEU-4 le plus faible (0,080), malgr\'e une adh\'erence
parfaite au format (100\,\%)}. Une investigation approfondie de la base de connaissances
en explique l'origine~:
\begin{itemize}[leftmargin=*]
  \item Sur les 261 entr\'ees PostgreSQL, seules 38 entr\'ees disposent d'une description
        technique riche r\'edig\'ee manuellement~; les 223 autres sont issues directement
        du libell\'e laconique du fichier \texttt{errcodes.txt} (\og{}\emph{generic template}\fg{}).
  \item Lorsque le r\'ecup\'erateur BM25 injecte un contexte pauvre, le mod\`ele T5
        g\'en\`ere des recommandations diagnostiques g\'en\'eralistes valides, mais dont le
        vocabulaire diff\`ere de la cible de test, p\'enalisant lourdement le score BLEU-4.
  \item \`A l'inverse, MySQL (0,170) b\'en\'eficie de 79 codes int\'egralement cur\'es avec
        des causes et solutions tr\`es pr\'ecises, permettant un meilleur alignement lexical.
\end{itemize}

\section{D\'emarche it\'erative de correction}

\subsection{D\'etection et correction des biais de r\'ecup\'eration inter-moteurs}

BM25 sans filtrage renvoyait des entr\'ees Oracle pour des logs PostgreSQL
(mots communs \og{}archive\fg{}, \og{}log\fg{}). Correction~: filtrage strict
par moteur dans \texttt{retrieve()}.

\subsection{D\'etection et correction des erreurs de formatage temporel (MySQL)}

Le parseur MySQL produisait des horodatages incorrects pour les logs h\'erit\'es.
Correction~: ajout d'un second pattern regex pour les deux formats de timestamp.

\subsection{D\'etection et correction du regroupement incorrect des pseudo-codes}

Les pseudo-codes \texttt{PG-ERROR} \'etaient regroup\'es sous une m\^eme cl\'e,
perdant le contexte sp\'ecifique. Correction~: cl\'e composite (code + message).

\subsection{D\'etection et correction de la fuite de valeurs de substitution}

Des placeholders non r\'esolus (\texttt{\%s}) dans les donn\'ees d'entra\^inement
\'etaient reproduits dans les sorties T5. Correction~: r\'esolution \`a la
g\'en\'eration des donn\'ees + post-traitement de nettoyage.

\subsection{Synth\`ese des anomalies identifi\'ees et corrig\'ees}

L'Annexe~C pr\'esente le d\'etail chronologique complet de toutes les corrections.

\section{Analyse qualitative des rapports g\'en\'er\'es}
\label{sec:qualitatif}

\subsection{Cas de succ\`es~: correspondances exactes}

Pour les codes bien couverts (ORA-01555, ORA-00060, ORA-00257), les explications
sont pr\'ecises, structur\'ees et directement actionnables~:

\begin{lstlisting}[language={},caption={Extrait de rapport pour ORA-01555},label=lst:report_example]
MEANING: snapshot too old: rollback segment number 4 with name "RBS3"
LIKELY_CAUSE: Transaction longue lisant un snapshot avant recyclage du segment.
SUGGESTED_SOLUTION: Augmenter UNDO_RETENTION, agrandir les tablespaces UNDO.
CONFIDENCE: high
Source: kb_exact_match | KB refs: ORA-01555
\end{lstlisting}

\subsection{Cas d'\'echec r\'esiduels et leur analyse}

Pour les codes absents de la KB, le syst\`eme adopte un comportement conservateur~:
correspondance BM25 la plus proche signal\'ee comme \og{}suggestion non confirm\'ee\fg{},
confiance abaiss\'ee \`a \texttt{low}. Ce comportement est pr\'ef\'erable \`a une
explication invent\'ee.

\subsection{Benchmark d'hallucination et comportement face \`a des codes inconnus}

Un benchmark quantitatif d'abstention a \'et\'e conduit sur un \'echantillon de
\textbf{20 codes d'erreur in\'edits et synth\'etiquement perturb\'es} inject\'es dans le pipeline~:
\begin{itemize}[leftmargin=*]
  \item \textbf{0 / 20 hallucination av\'er\'ee}~: aucune fabrication de faux num\'ero d'erreur
        interne ni d'affirmation p\'eremptoire infond\'ee n'a \'et\'e observ\'ee~;
  \item \textbf{17 / 20 r\'eponses prudentes calibr\'ees (\emph{hedged})}~: le mod\`ele a
        g\'en\'er\'e un diagnostic g\'en\'erique assorti d'un niveau de confiance r\'eduit
        (\texttt{CONFIDENCE: low} ou \texttt{medium}) et du label de source
        \texttt{t5\_local\_no\_kb\_match}~;
  \item \textbf{3 / 20 replis conservateurs standards}~: recommandations d'inspection
        g\'en\'erale des traces syst\`eme sans sp\'eculation.
\end{itemize}

\section{Limites identifi\'ees et discussion honn\^ete}
\label{sec:limites}

\subsection{D\'es\'equilibre de couverture et profondeur des bases de connaissances}

Oracle dispose de 27\,282 entr\'ees exhaustives, MySQL de 79 entr\'ees cur\'ees, et
PostgreSQL de 261 entr\'ees (dont seulement 38 r\'eellement enrichies). Comme
d\'emontr\'e en section~5.2.3, cette disparit\'e de profondeur documentaire explique
directement le faible score BLEU-4 de PostgreSQL (0,080) par rapport \`a Oracle (0,210)
et MySQL (0,170). L'enrichissement manuel des 223 codes SQLSTATE g\'en\'eriques de
PostgreSQL constitue la priorit\'e imm\'ediate.

\subsection{D\'ependance \`a des donn\'ees d'entra\^inement synth\'etiques}

Les logs de production pr\'esentent probablement des caract\'eristiques non captur\'ees
dans le corpus synth\'etique (codes tr\`es rares, formulations atypiques).

\subsection{Absence de validation sur des logs de production r\'eels}

La validation sur des logs r\'eels est la seule mani\`ere de confirmer l'utilit\'e
op\'erationnelle r\'eelle du syst\`eme. Non r\'ealisable dans le cadre de ce stage
pour des raisons de confidentialit\'e.

\subsection{Risques r\'esiduels de g\'en\'eration incorrecte}

Pour les correspondances BM25 uniquement, le mod\`ele T5 peut g\'en\'erer des
explications plausibles mais incorrectes. Le m\'ecanisme de confiance att\'enue ce
risque mais ne l'\'elimine pas.

% ============================================================
% CONCLUSION
% ============================================================
\chapter*{Conclusion g\'en\'erale et perspectives}
\addcontentsline{toc}{chapter}{Conclusion g\'en\'erale et perspectives}

\section*{Synth\`ese des objectifs atteints}

Ce stage a permis de mener de bout en bout la conception et le d\'eveloppement
d'une plateforme compl\`ete de diagnostic automatique des logs de bases de donn\'ees
h\'et\'erog\`enes. Les objectifs fix\'es ont tous \'et\'e atteints~:
\begin{itemize}[leftmargin=*]
  \item Support multi-moteurs (Oracle, PostgreSQL, MySQL) avec d\'etection automatique~;
  \item Base de connaissances de 27\,622 entr\'ees sans donn\'ees de production~;
  \item Mod\`ele FLAN-T5 fine-tun\'e, autonome, sans API externe payante~;
  \item Interface web d\'eploy\'ee en production sur Railway~;
  \item 55~tests automatis\'es, 100\,\% de r\'eussite.
\end{itemize}

Sur le plan technique, les principaux apports sont la ma\^itrise du cycle complet
d'un syst\`eme NLP appliqu\'e~: corpus, fine-tuning, \'evaluation, correction de
bugs subtils et d\'eploiement en production.

\section*{Perspectives d'am\'elioration}

\begin{itemize}[leftmargin=*]
  \item \textbf{Extension de la couverture MySQL et PostgreSQL}~: automatiser
        l'extraction des codes MySQL comme pour Oracle~;
  \item \textbf{Retrieval hybride}~: combiner BM25 et embeddings denses pour am\'eliorer
        la qualit\'e de r\'ecup\'eration pour les codes absents~;
  \item \textbf{Corr\'elation d'\'ev\'enements}~: relier plusieurs codes co-occurring
        pour identifier une cause racine commune~;
  \item \textbf{Scoring de s\'ev\'erit\'e dynamique}~: score global pour l'ensemble
        du log analys\'e~;
  \item \textbf{Validation sur logs r\'eels}~: priorit\'e absolue d\`es qu'un
        protocole de pseudonymisation sera valid\'e avec la DCSI.
\end{itemize}

% ============================================================
% BIBLIOGRAPHIE
% ============================================================
\begin{thebibliography}{99}

\bibitem{oracledoc}
Oracle Corporation.
\textit{Oracle Database Error Messages}.
Oracle Help Center, 2024.
\url{https://docs.oracle.com/en/database/oracle/oracle-database/}

\bibitem{bs4}
Richardson, L.
\textit{Beautiful Soup Documentation}.
Crummy.com, 2023.
\url{https://pypi.org/project/beautifulsoup4/}

\bibitem{bm25}
Brown, D.
\textit{rank\_bm25~: A Collection of BM25 Algorithms in Python}.
GitHub, 2020.
\url{https://github.com/dorianbrown/rank_bm25}

\bibitem{bm25orig}
Robertson, S. E. et Zaragoza, H.
\textit{The Probabilistic Relevance Framework: BM25 and Beyond}.
Foundations and Trends in Information Retrieval, 3(4), 333--389, 2009.

\bibitem{lewis2020rag}
Lewis, P., Perez, E., Piktus, A., et al.
\textit{Retrieval-Augmented Generation for Knowledge-Intensive NLP Tasks}.
Advances in Neural Information Processing Systems (NeurIPS), 33, 9459--9474, 2020.

\bibitem{zhu2019log}
Zhu, J., He, S., Liu, J., He, P., Xie, Q., Zheng, Z. et Lyu, M. R.
\textit{Tools and Benchmarks for Automated Log Parsing}.
Proceedings of the 41st IEEE/ACM International Conference on Software Engineering (ICSE), 121--130, 2019.

\bibitem{le2022log}
Le, V.-H. et Zhang, H.
\textit{Log-based Anomaly Detection with Deep Learning: How Far Are We?}
Proceedings of the 44th IEEE/ACM International Conference on Software Engineering (ICSE), 1356--1367, 2022.

\bibitem{flanT5}
Chung, H. W. et al.
\textit{Scaling Instruction-Finetuned Language Models}.
Journal of Machine Learning Research (JMLR), 2024.

\bibitem{fastapi}
Ram\'irez, S.
\textit{FastAPI Documentation}.
Tiangolo, 2024.
\url{https://fastapi.tiangolo.com/}

\bibitem{hfhub}
Wolf, T. et al.
\textit{Transformers: State-of-the-Art Natural Language Processing}.
Proceedings of EMNLP (Systems Demonstrations), 38--45, 2020.

\end{thebibliography}

% ============================================================
% ANNEXES
% ============================================================
\appendix

\chapter{Extraits de code significatifs}

\section{Parseur Oracle --- extraction des codes d'erreur}

\begin{lstlisting}[language=Python,caption={Extrait de src/log\_parser.py},label=lst:parser]
# 81 prefixes reconnus : ORA, TNS, RMAN, PLS, CRS, DRG, ...
ORACLE_PREFIXES = [
    "ORA", "TNS", "RMAN", "PLS", "CRS", "DRG", "INS",
    "OCI", "NZE", "QSM", "PRVT", ...
]
_ORACLE_PATTERN = re.compile(
    r'\b(' + '|'.join(ORACLE_PREFIXES) + r')-(\d{3,6})\b'
)

def extract_oracle_errors(text: str) -> list[ErrorOccurrence]:
    results = []
    for m in _ORACLE_PATTERN.finditer(text):
        code = normalize_code(m.group(0))
        results.append(ErrorOccurrence(code=code, ...))
    return deduplicate(results)
\end{lstlisting}

\section{Pipeline de g\'en\'eration --- fonction \texttt{generate\_t5()}}

\begin{lstlisting}[language=Python,caption={Extrait de src/rag/generator.py},label=lst:generate_t5]
def generate_t5(code, raw_line, context, retrieved,
                model_dir=DEFAULT_T5_MODEL_DIR,
                max_new_tokens=80, engine="oracle"):
    instruction = INSTRUCTIONS.get(engine, INSTRUCTIONS["oracle"])
    try:
        tokenizer, model = _load_t5(model_dir)
    except Exception:
        fallback = generate_mock(code, raw_line, context, retrieved)
        fallback.source = "t5_fallback_to_kb"
        return fallback

    kb_lines = [
        f"- [{e.code}] {e.message}\n"
        f"  cause: {e.cause}\n"
        f"  solution: {e.solution}\n"
        f"  (score: {s:.1f})"
        for e, s in (retrieved or [])
    ]
    kb_block = "\n".join(kb_lines) or "(no entries found)"
    source_text = instruction + "\n\n" + context + "\n\n" \
                  + "RETRIEVED KNOWLEDGE:\n" + kb_block

    inputs = tokenizer(source_text, return_tensors="pt",
                       truncation=True, max_length=512).to(model.device)
    import torch
    with torch.no_grad():
        output_ids = model.generate(
            **inputs, max_new_tokens=max_new_tokens,
            do_sample=False, repetition_penalty=1.2,
            no_repeat_ngram_size=3,
        )
    text = tokenizer.decode(output_ids[0], skip_special_tokens=True)
    return _parse_and_build_explanation(text, code, retrieved)
\end{lstlisting}

\chapter{Exemples de rapports g\'en\'er\'es par moteur}

\section{Rapport Oracle --- ORA-00060 et ORA-01555}

\begin{lstlisting}[language={},caption={Extrait de rapport Oracle},label=lst:report_oracle]
## ORA-00060  (seen 3x, lines: [142, 287, 301])
MEANING: Deadlock detected while waiting for resource.
LIKELY_CAUSE: Deux transactions ont tente d'acquerir des verrous croises.
SUGGESTED_SOLUTION: Revoir l'ordre d'acquisition des verrous. Examiner
  V$LOCK et V$SESSION pour identifier les sessions impliquees.
CONFIDENCE: high | Source: kb_exact_match | KB refs: ORA-00060

## ORA-01555  (seen 7x, lines: [89, 112, ...])
MEANING: snapshot too old: rollback segment number 4 with name "RBS3"
LIKELY_CAUSE: Requete longue lisant un snapshot avant recyclage du segment.
SUGGESTED_SOLUTION: Augmenter UNDO_RETENTION, agrandir les tablespaces UNDO.
CONFIDENCE: high | Source: kb_exact_match | KB refs: ORA-01555
\end{lstlisting}

\section{Rapport PostgreSQL --- 42P01 et 53300}

\begin{lstlisting}[language={},caption={Extrait de rapport PostgreSQL},label=lst:report_pg]
## 42P01  (seen 2x, lines: [34, 78])
MEANING: undefined_table: La relation specifiee n'existe pas.
LIKELY_CAUSE: Nom de table incorrect ou table non creee dans le schema courant.
SUGGESTED_SOLUTION: Verifier le nom de la table via \dt ou pg_tables.
CONFIDENCE: high | Source: kb_exact_match

## 53300  (seen 1x, lines: [189])
MEANING: too_many_connections: Limite de connexions atteinte.
SUGGESTED_SOLUTION: Augmenter max_connections dans postgresql.conf
  ou deployer PgBouncer comme pooler de connexions.
CONFIDENCE: high | Source: kb_exact_match
\end{lstlisting}

\chapter{D\'etail chronologique des corrections apportées au pipeline}

\begin{enumerate}[leftmargin=*]
  \item \textbf{Bug parseur regex}~: liste des pr\'efixes Oracle pass\'ee de 13 \`a 81.
  \item \textbf{Bug tokeniseur T5}~: sauts de ligne supprim\'es bloquant le parseur
        de r\'eponse. Corrig\'e par regex \emph{lookahead}.
  \item \textbf{Bug contamination inter-moteurs}~: BM25 sans filtrage. Corrig\'e
        par filtrage dans \texttt{retrieve()}.
  \item \textbf{Bug pseudo-codes PG-ERROR}~: regroupement sous m\^eme cl\'e.
        Corrig\'e par cl\'e composite.
  \item \textbf{Bug placeholders non r\'esolus}~: valeurs \texttt{\%s} propag\'ees
        dans les sorties T5. Corrig\'e par r\'esolution \`a la g\'en\'eration.
  \item \textbf{Bug reload=True}~: file watcher Uvicorn red\'emarrait le serveur.
        Corrig\'e par \texttt{reload} conditionnel.
  \item \textbf{Bug KB rebuild par requ\^ete}~: +2 \`a 4~s de latence.
        Corrig\'e par singletons au d\'emarrage.
  \item \textbf{Bug OOM torch.compile()}~: +300--500~Mo de m\'emoire.
        Corrig\'e par suppression de \texttt{torch.compile()}.
\end{enumerate}

\chapter{Guide d'installation et de d\'eploiement}

\section{Installation locale}

\begin{lstlisting}[language=bash,caption={Installation locale},label=lst:install]
# 1. Cloner le depot
git clone https://github.com/rayenthabet004-spec/Projet-TT.git
cd Projet-TT

# 2. Installer les dependances Python
pip install -r requirements.txt

# 3. Lancer le serveur de developpement
ENV=development python web_app.py
# Acces : http://localhost:8000

# 4. Lancer la suite de tests
python -m pytest tests/ -v
# Attendu : 55 passed
\end{lstlisting}

\section{D\'eploiement Docker}

\begin{lstlisting}[language=bash,caption={Docker build et run},label=lst:docker]
# Build (pre-cache le modele T5 dans l'image)
docker build -t projet-tt .

# Lancement
docker run -p 8000:8000 -e HF_TOKEN=<votre_token> projet-tt
\end{lstlisting}

\section{Variables d'environnement}

\begin{table}[H]
  \centering
  \caption{Variables d'environnement de l'application}
  \label{tab:env_vars}
  \begin{tabularx}{\textwidth}{llX}
    \toprule
    \textbf{Variable} & \textbf{D\'efaut} & \textbf{Description} \\
    \midrule
    \texttt{PORT}             & 8000       & Port d'\'ecoute \\
    \texttt{ENV}              & production & \texttt{development} active le rechargement \\
    \texttt{HF\_HOME}         & /app/hf\_cache & Cache Hugging Face \\
    \texttt{HF\_HUB\_OFFLINE} & 1          & D\'esactive les appels r\'eseau HF Hub \\
    \texttt{HF\_TOKEN}        & (vide)     & Token HF (requis au build si d\'ep\^ot priv\'e) \\
    \texttt{T5\_HF\_REPO\_ID} & rayenthabet004/tt-multi-engine-t5 & Identifiant mod\`ele \\
    \texttt{GEMINI\_API\_KEY} & (vide)     & Cl\'e Gemini (optionnel) \\
    \texttt{GROQ\_API\_KEY}   & (vide)     & Cl\'e Groq (optionnel) \\
    \bottomrule
  \end{tabularx}
\end{table}

\end{document}
